% BHU PYQ -- Manually transcribed & error-corrected
% Paper: BCF-313 : Mutual Fund Operations
% Exam: B.Com. (Hons.) FMM Semester V Examination, 2022-23

\documentclass[11pt,a4paper]{article}
\usepackage[a4paper,top=2.5cm,bottom=2.5cm,left=2.8cm,right=2.8cm,headheight=14pt]{geometry}
\usepackage{amsmath,amssymb}
\usepackage[T1]{fontenc}
\usepackage[utf8]{inputenc}
\usepackage{lmodern,microtype,fancyhdr,enumitem,hyperref,lastpage,parskip}

\hypersetup{
    colorlinks=true,
    urlcolor=black,
    linkcolor=black,
    pdftitle={BCF-313: Mutual Fund Operations -- BHU B.Com. (Hons.) FMM Sem V 2022-23}
}

\pagestyle{fancy}
\fancyhf{}
\renewcommand{\headrulewidth}{0.4pt}
\renewcommand{\footrulewidth}{0.4pt}
\fancyhead[L]{\small\textit{Banaras Hindu University}}
\fancyhead[R]{\small\textit{BCF-313: Mutual Fund Operations}}
\fancyfoot[C]{\small Page \thepage\ of \pageref{LastPage}}

\newlist{parts}{enumerate}{2}
\setlist[parts,1]{label=(\alph*),leftmargin=*,itemsep=5pt,topsep=3pt}
\setlist[parts,2]{label=(\roman*),leftmargin=*,itemsep=3pt,topsep=2pt}
\newcommand{\pts}[1]{\hfill\textit{[#1]}}

\begin{document}

\begin{center}
  {\large\bfseries BANARAS HINDU UNIVERSITY}\\[2pt]
  {\normalsize B.Com. (Hons.) FMM --- Semester V Examination, 2022-23}\\[6pt]
  \rule{0.8\linewidth}{0.5pt}\\[6pt]
  {\large\bfseries Paper No. BCF-313: Mutual Fund Operations}\\[4pt]
  {\normalsize Time: 3 Hours\hfill Maximum Marks: 70}
\end{center}

\rule{\linewidth}{0.4pt}

\smallskip
\noindent\textit{Instructions: Attempt all questions. All questions carry equal marks (14 marks each).}

\bigskip

\noindent\textbf{Question 1.}
Define mutual funds. Also describe the benefits of mutual funds in the current Indian financial system.\pts{14}

\centerline{\textbf{OR}}

Differentiate between the two:\pts{14}
\begin{parts}
  \item Growth fund and Index fund
  \item Public sector funds and private sector funds
\end{parts}

\medskip\rule{\linewidth}{0.2pt}

\noindent\textbf{Question 2.}
Briefly explain the risk and return analysis of mutual fund schemes while keeping in mind the volatile financial environment.\pts{14}

\centerline{\textbf{OR}}

Write the important provisions of the Indian Trusts Act taking consideration of mutual funds.\pts{14}

\medskip\rule{\linewidth}{0.2pt}

\noindent\textbf{Question 3.}
Explain and illustrate the calculations of Net Asset Value (NAV) in mutual funds.\pts{14}

\centerline{\textbf{OR}}

What are Asset Management Companies (AMC)? What are their functions and obligations in respect of mutual fund operations?\pts{14}

\medskip\rule{\linewidth}{0.2pt}

\noindent\textbf{Question 4.}
Explain the basic structure of mutual fund management. Also mention the various participants of mutual fund operations.\pts{14}

\centerline{\textbf{OR}}

Write short notes on any two of the following:\pts{14}
\begin{parts}
  \item Accounting policies and standards
  \item Inspection and audit of mutual funds
  \item Procedure for action in case of default
\end{parts}

\medskip\rule{\linewidth}{0.2pt}

\noindent\textbf{Question 5.}
``Mutual funds are the sharpest weapons in the armory of a financial planner.'' Explain the statement.\pts{14}

\centerline{\textbf{OR}}

What is 'Risk'? What are the various kinds of risks prevalent on mutual fund investments?\dots\pts{14}

\vfill
\rule{\linewidth}{0.4pt}
\begin{center}\small
\href{https://bhu-exam.vercel.app/}{Click here to visit our website}\\[4pt]
\href{https://me-aryan.vercel.app/}{Click here to visit our contributor}\\[4pt]
\href{https://quashan-me.vercel.app/}{Click here to visit our contributor}
\end{center}

\end{document}
